本研究においては、先行研究で報告されている旗にモデルが過学習してしまう問題を解決するために異なる物体検出モデルを用いて３つの検証を行った。
\subsection{方法1:検出性能検証}
 まず、本研究において導入した物体検出モデルの基本的な検出性能を評価するために、実際に圃場のデータセットで学習を行い、病害木の検出を行った。
\subsection{方法2:旗問題検証}
次に、
\subsection{方法3:HBBT検証}
本研究において使用したドローン空撮画像のデータセットではすべての感染木に目印として旗が設置されている。これによって、感染木と人工物が常に共起する。これによって物体検出モデルが病害に起因しない視覚的特徴である旗を病害木の識別手掛かりとして過学習してしまう現象を誘発してしまう。そのような旗問題に対し、先行研究とは異なるモデルを導入しその挙動を確認することでモデル自体の性能により旗問題を解決できるか検証する。



モデルの図やポンチ絵を書くときは google 図形描画\footnote{\url{https://docs.google.com/drawings/create?usp=apps_start&hl=ja}} や diagram.net\footnote{\url{https://www.diagrams.net/}} を使って図を書いて PDF にすることを解像度の高い図を埋め込めます。

\subsection


\subsubsection{方法2の詳細1：旗問題}


\subsubsection{方法2の詳細2}

方法3の詳細について記述する．
この方法に対してこのような実験を行いました．

\input{images/tex/dataset}
